\section{Introduction}
\label{sec:intro}

The emergence of massive pre-trained foundational models has transformed the paradigm of modern machine learning. In practice, practitioners fine-tune these base models on various downstream tasks, producing specialized expert models. While highly performant on their respective target domains, deployment of multiple separate experts incurs unsustainable computational and storage footprints. To circumvent this, \emph{model merging} has emerged as an exceptionally elegant and cost-effective paradigm, consolidating multiple task-specific expert weights into a single, unified multi-task model without requiring additional training or paired joint optimization data \cite{ilharco2022editing, wortsman2022robust, choshen2022fusing}.

Among various merging strategies, \emph{Task Arithmetic} \cite{ilharco2022editing} has established a powerful baseline by representing task-specific experts as directions in parameter space---termed \emph{task vectors}---and linearly adding them to the pre-trained base weights. Formally, for a base model parameterized by $\Theta_{\text{base}}$ and task-specific experts $\Theta_k$ for $k \in \{1, \dots, K\}$, the consolidated model parameters are obtained via:
\begin{equation}
\Theta_{\text{merged}} = \Theta_{\text{base}} + \sum_{k=1}^K \lambda_k (\Theta_k - \Theta_{\text{base}})
\end{equation}
where $\lambda_k \ge 0$ represents the merging coefficient for task $k$. Despite its simplicity and empirical success, static and uniform merging (where $\lambda_k$ is constant across all layers) overlooks a fundamental property of deep neural networks: different layers exhibit vastly different sensitivities, functional behaviors, and representation capacities \cite{raghu2017svcca}. Consequently, uniform scaling forces a sub-optimal trade-off between conflicting tasks.

To address this limitation, recent works have proposed dynamic, layer-wise coefficient optimization at test time. For instance, \emph{AdaMerging} \cite{lu2023adamerging} optimizes separate merging coefficients $\lambda_{k, l}$ for each task $k$ and layer $l$ using unlabeled test-time adaptation (TTA) streams. By minimizing a self-supervised surrogate objective---typically the Shannon entropy of model predictions---over incoming test batches, AdaMerging dynamically aligns the model to the test distribution. 

However, we expose an on-the-fly, unconstrained layer-wise merging overfitting phenomenon, which we term the \textbf{Overfitting-Optimizer Paradox}. Because test-time adaptation streams are processed in small, localized batches, the optimization objective is inherently transductive, corrupted by high-frequency local sampling noise. Unconstrained optimizers possessing a large number of independent degrees of freedom (e.g., $K \times L$ free parameters) quickly exploit this transductive noise to minimize the surrogate entropy objective. Instead of aligning to the underlying semantic sensitivity profile of the network, the unconstrained optimizer memorizes high-frequency noise, causing the merging coefficients to oscillate wildly. This leads to what we define as \emph{representation collapse}---a severe degradation in the model's fundamental generalization capability, particularly on challenging domains where the accuracy can drop below the non-adaptive baseline.

To mitigate this overfitting, researchers have attempted to enforce structural smoothness on the merging coefficients. One notable approach is \emph{PolyMerge} \cite{polymerge2024}, which constrains layer-wise coefficients to a low-dimensional continuous subspace parameterized as low-degree polynomials of normalized depth. While this restriction successfully filters out high-frequency transductive noise, it introduces a severe, previously unaddressed numerical vulnerability: \emph{exponential ill-conditioning}. In PolyMerge, the mapping from spectral parameters to spatial coefficients is mediated by a standard monomial basis ($1, \bar{l}, \bar{l}^2, \dots$), yielding a Vandermonde-type design matrix. The condition number of the corresponding Gram matrix scales exponentially as $\mathcal{O}(4^d)$ where $d$ is the polynomial degree. For a cubic polynomial ($d=3$), the condition number exceeds $10,400$. This extreme ill-conditioning compresses the optimization landscape, creating highly anisotropic "stiff" loss valleys that destabilize gradient descent, trigger vanishing or exploding gradients, and severely compromise convergence speed and reliability.

In this work, we propose \textbf{ChebyMerge} (Stable and Optimal Continuous Subspace Model Merging) to resolve both the Overfitting-Optimizer Paradox and the numerical ill-conditioning of monomial parameterizations. Grounded in orthogonal approximation theory, ChebyMerge projects the layer-wise model-merging coefficients onto an orthogonal subspace spanned by Chebyshev polynomials of the first kind ($T_j(x)$). 

By framing model merging as a continuous spectral approximation problem under a Chebyshev basis, we unlock three profound advantages:
\begin{enumerate}
    \item \textbf{Minimax Optimality:} Chebyshev polynomials provide near-optimal uniform approximation under the supremum norm ($L_\infty$), guaranteeing that our low-dimensional parameterization minimizes the maximum possible representation error for any smooth underlying layer-wise importance profile.
    \item \textbf{Perfect Numerical Conditioning:} Because Chebyshev polynomials are orthogonal under the continuous weighted inner product, their evaluations on a discrete grid remain nearly orthogonal. The condition number of the Chebyshev Gram matrix is bounded by a tiny constant close to 1 (e.g., $\approx 2.95$ for $d=3$), achieving a highly significant \textbf{3,527$\times$ improvement} over PolyMerge's monomial basis. This ensures an isotropic optimization landscape and exceptionally stable gradient-based convergence.
    \item \textbf{Implicit Boundary Sensitivity Matching:} The roots and extrema of Chebyshev polynomials naturally cluster near the boundaries of the interval $[-1, 1]$. In deep neural networks, early and deep layers are highly sensitive to parameter modifications, whereas intermediate layers are relatively flat \cite{zhang2019all}. ChebyMerge's mathematical structure naturally concentrates representational resolution at these highly sensitive boundaries, matching the intrinsic physical sensitivity profile of deep models.
\end{enumerate}

To evaluate these methods under rigorous, mathematically controlled conditions, we present a physically-grounded coupled non-convex loss-landscape and optimization simulator. Rather than evaluating on black-box networks where internal gradient fields are unobservable and ground-truth sensitivities are unknown, our simulator emulates the exact layer sensitivity profiles, non-convex Rastrigin-type loss surfaces, inter-layer covariance couplings, and multi-scale transductive noise characteristics of deep Vision Transformers. This provides an essential methodological tool for isolating optimization and numerical conditioning dynamics under perfect ground-truth control.

Our empirical results and subsequent theoretical analysis reveal an insight regarding the tension between optimization conditioning and regularization, which we refer to as the \textbf{Conditioning-Generalization Paradox}. We expose that PolyMerge's extreme ill-conditioning acts as an accidental, uncontrolled \emph{implicit spectral damping filter}, which prevents the optimizer from adjusting high-frequency terms, thereby acting as a beneficial regularizer against transductive noise. However, relying on numerical instability for regularization is unprincipled, fragile, and sensitive to hyperparameter choices. In contrast, ChebyMerge decouples numerical conditioning from regularization: it ensures perfectly isotropic and stable gradient flow, enabling researchers to apply \emph{explicit, controllable regularization} (such as spectral weight decay or lower polynomial degrees) rather than relying on accidental numerical errors. 

Our core findings demonstrate that:
\begin{itemize}
    \item Unconstrained TTA optimizers suffer severely from the Overfitting-Optimizer Paradox, collapsing to an average generalization score of $78.67\%$ under simulated non-convex stress tests (with SVHN collapsing to $55.30\%$).
    \item Continuous subspace projections completely prevent this collapse, with ChebyMerge ($d=2$) achieving a robust generalization score of $85.25\%$, outperforming uniform Task Arithmetic ($84.44\%$) and unconstrained Adam ($78.67\%$).
    \item The well-conditioned optimization landscape of ChebyMerge enables smooth, stable, and rapid convergence, establishing it as a mathematically rigorous and highly practical method for test-time model merging.
\end{itemize}
